\subsection{DEFINITION OF AFFINE SPACES}

DEFINITION 1. Given a left (resp. right) vector space T over a field K, an affine space attached to T is any homogeneous space E of the additive group T (I, §5, no. 5) such that 0 is the only operator in T leaving invariant all the elements of E (that is, T operates faithfully and transitively on E). Under these conditions, T is called the translation space of E and its elements are called the translations of E (or free vectors of E).

In what follows we shall confine our attention to the case where T is a left vector space over K. The dimension (over K) of the translation vector space T of an affine space E is called the dimension of E (over K) and is denoted by dim E or dim \(_{K}\) E. An affine space of dimension one (resp. two) is called an affine line (resp. an affine plane). The elements of an affine space are also called points.

Under the conditions of Definition 1, for \(\mathbf{t} \in \mathbf{T}\) and \(a \in \mathbf{E}\) we shall denote by \(\mathbf{t} + a\) or \(a + \mathbf{t}\) the transform of the point \(a\) under \(\mathbf{t}\) . Then the relations

\[
\mathbf {s} + (\mathbf {t} + a) = (\mathbf {s} + \mathbf {t}) + a, \quad 0 + a = a\tag{1}
\]

hold for \(\mathbf{s} \in \mathrm{T}\) , \(\mathbf{t} \in \mathrm{T}\) , \(a \in \mathrm{E}\) . The mapping \(x \mapsto x + \mathbf{t}\) is a bijection of \(\mathbf{E}\) onto itself, which we identify with \(\mathbf{t}\) . Definition 1 moreover implies that, for all \(a \in \mathrm{E}\) , the mapping \(\mathbf{t} \mapsto \mathbf{t} + a\) is a bijection of \(\mathrm{T}\) onto \(\mathbf{E}\) . In other words, given two points \(a\) , \(b\) of \(\mathbf{E}\) , there exists one and only one translation \(\mathbf{t}\) such that \(b = \mathbf{t} + a\) ; we shall denote this translation by \(b - a\) ; then the formulae

\[
a - a = 0, \quad a - b = - (b - a), \quad b = (b - a) + a\tag{2}
\]

\[
(c - b) + (b - a) = c - a
\]

hold for \(a \in \mathbf{E}\) , \(b \in \mathbf{E}\) , \(c \in \mathbf{E}\) . If four points \(a, b, a', b'\) of \(\mathbf{E}\) are such that \(b - a = b' - a'\) , the formula

\[
b ^ {\prime} = \left(b ^ {\prime} - b\right) + (b - a) + a = \left(b ^ {\prime} - a ^ {\prime}\right) + \left(a ^ {\prime} - a\right) + a
\]

and the commutativity of addition in T show that \(b' - b = a' - a\) .

Given a point \(a \in E\) , the mapping \(x \mapsto x - a\) is a bijection of \(E\) onto \(T\) ; when \(E\) is identified with \(T\) under this mapping, we say that \(E\) is considered as the vector space obtained by taking \(a\) as origin in \(E\) . Conversely, every vector space \(T\) has canonically the structure of an affine space attached to \(T\) , namely the homogeneous space structure corresponding to the subgroup \(\{0\}\) of \(T\) (I, § 5, no. 6).

Remark. The definitions of this no. and some of the results which follow extend immediately to the case where, instead of a vector space T, we consider an arbitrary commutative group with operators T.

